\section{Contributions}

We introduce a benchmark-construction pipeline for multilingual legal retrieval using the JRC-Acquis corpus. Our work makes five main contributions. First, we construct a document-level multilingual retrieval corpus from aligned European Union legal texts by treating \texttt{CELEX} as the cross-lingual document anchor and preserving one retrieval document per \((\texttt{CELEX}, \texttt{language})\) combination, with an optional EuroVoc-based chemical-domain filter for chemistry-focused legal experiments. Second, we develop a legal-text preprocessing procedure that removes non-substantive artifacts, separates body, annex, and signature material, and distinguishes between retrieval context and QA-generation context rather than exposing the model to a single undifferentiated text field. Third, we introduce a \texttt{CELEX}-group sampling and subset-construction strategy that selects source documents from configurable strict or soft QA-candidate pools, prepares oversampled replacement candidates, and builds a retained multilingual subset over which the final relevance labels are defined. Fourth, we formulate supervision through same-language legal question--answer generation from selected target-side document realizations, together with a staged validation stack that filters language errors, faithfulness failures, weak retrieval questions, undesirable legal question shapes, and translated documents that do not support the generated answer. Fifth, we package the resulting benchmark in standard retrieval form---\texttt{corpus}, \texttt{queries}, \texttt{qrels}, and \texttt{qac}---and evaluate it through a Hugging Face-backed MTEB retrieval setup with multilingual encoder baselines, variant-specific qrels, comparison reports, and a local numerical audit report for construction statistics.
